\section*{Parte 1: Sistema de misiones de PurisQuest}

\subsection*{{\small \faDragon}\; Contexto}

PurisQuest es un videojuego de rol ambientado en un mundo de fantasía medieval.
Cada jugador comienza como un \emph{Aprendiz} y avanza de nivel conforme completa misiones y acumula puntos de experiencia (XP).
En este examen debe implementar el núcleo del sistema de misiones del juego usando Python.

Debe entregar un archivo llamado \texttt{parte1\_NombreApellido.py} como solución.

\subsection*{{\small \faDatabase}\; Datos predefinidos}

El programa debe iniciar con el siguiente diccionario declarado directamente en el código:

\begin{minted}{python}
misiones = {
    "Bosque Maldito":   {"xp": 150, "dificultad": "Fácil"},
    "Cueva Oscura":     {"xp": 300, "dificultad": "Media"},
    "Torre Abandonada": {"xp": 500, "dificultad": "Difícil"},
    "Aldea Olvidada":   {"xp": 100, "dificultad": "Fácil"},
    "Minas Profundas":  {"xp": 250, "dificultad": "Media"},
    "Castillo Eterno":  {"xp": 800, "dificultad": "Legendaria"},
}
\end{minted}

Además, se deben declarar:

\begin{itemize}
    \item Una \textbf{lista vacía} \mintinline{python}{completadas} para registrar las misiones completadas.
    \item Una variable \mintinline{python}{xp_acumulado} inicializada en \mintinline{python}{0}.
\end{itemize}

El nivel del personaje se determina según la siguiente tabla:

\begin{table}[H]
    \centering
    \begin{tabular}{cc}
        \toprule
        \textbf{XP acumulado} & \textbf{Nivel} \\
        \midrule
        0 -- 299    & Aprendiz \\
        300 -- 699  & Guerrero \\
        700 -- 1199 & Héroe \\
        1200 o más  & Leyenda \\
        \bottomrule
    \end{tabular}
    \caption{Niveles del personaje según XP acumulado.}
    \label{tab:niveles}
\end{table}

\subsection*{{\small \faCogs}\; Funciones requeridas}

\subsubsection*{Función \texttt{calcular\_nivel(xp)}}

Recibe el XP acumulado y \textbf{retorna} el nombre del nivel correspondiente según la tabla anterior:
\mintinline{python}{"Aprendiz"}, \mintinline{python}{"Guerrero"}, \mintinline{python}{"Héroe"} o \mintinline{python}{"Leyenda"}.

\subsubsection*{Función \texttt{mostrar\_misiones(misiones)}}

Recibe el diccionario de misiones y recorre sus elementos para imprimir el nombre, el XP y la dificultad de cada misión.

\subsubsection*{Función \texttt{completar\_mision(nombre, misiones, completadas, xp)}}

Recibe el nombre ingresado por el usuario, el diccionario de misiones, la lista de completadas y el XP actual.
Aplica la siguiente lógica en orden:

\begin{enumerate}
    \item Normalizar el nombre para aceptar cualquier combinación de mayúsculas y minúsculas.
    \item Si \textbf{no existe} en el diccionario, imprimir y no hacer nada más:
    \begin{minted}{text}
Misión no encontrada.
    \end{minted}
    \item Si \textbf{ya está} en la lista \mintinline{python}{completadas}, imprimir y no hacer nada más:
    \begin{minted}{text}
Esa misión ya fue completada.
    \end{minted}
    \item Si la misión es válida, agregar el nombre a \mintinline{python}{completadas}, sumar el XP de esa misión al XP recibido e imprimir:
    \begin{minted}{text}
Misión completada. +<XP> XP
    \end{minted}
    \item \textbf{Retornar} el nuevo valor de XP acumulado en todos los casos.
\end{enumerate}

\subsubsection*{Función \texttt{mostrar\_perfil(jugador, xp, completadas)}}

Recibe el nombre del jugador, el XP acumulado y la lista de completadas.

\begin{minted}{text}
=== Perfil de <jugador> ===
XP total: <xp>
Nivel: <nivel>
Misiones completadas: <lista>
\end{minted}

\subsection*{{\small \faList}\; Menú principal}

Al iniciar el programa se solicita el nombre del jugador y se muestra el siguiente menú en un ciclo \mintinline{python}{while} hasta que el usuario elija salir:

\begin{minted}{text}
=== PurisQuest ===
1. Ver misiones disponibles
2. Completar una misión
3. Ver perfil del personaje
4. Salir
\end{minted}

Cada opción debe solicitar los datos necesarios e invocar la función correspondiente.
Al salir (opción 4), imprimir:

\begin{minted}{text}
Hasta la próxima, <jugador>. Nivel alcanzado: <nivel>.
\end{minted}

\subsection*{{\small \faPlay}\; Ejemplo de ejecución}

Los valores en \textbf{negrita} corresponden a lo que escribe el usuario; el resto es salida del programa.

\begin{minted}[escapeinside=||]{text}
Nombre del personaje: |\textbf{Wilson}|

=== PurisQuest ===
1. Ver misiones disponibles
2. Completar una misión
3. Ver perfil del personaje
4. Salir
Opción: |\textbf{1}|

Bosque Maldito - XP: 150 - Dificultad: Fácil
Cueva Oscura - XP: 300 - Dificultad: Media
Torre Abandonada - XP: 500 - Dificultad: Difícil
Aldea Olvidada - XP: 100 - Dificultad: Fácil
Minas Profundas - XP: 250 - Dificultad: Media
Castillo Eterno - XP: 800 - Dificultad: Legendaria

=== PurisQuest ===
...
Opción: |\textbf{2}|
Nombre de la misión: |\textbf{bosque maldito}|
Misión completada. +150 XP

=== PurisQuest ===
...
Opción: |\textbf{2}|
Nombre de la misión: |\textbf{Bosque Maldito}|
Esa misión ya fue completada.

=== PurisQuest ===
...
Opción: |\textbf{2}|
Nombre de la misión: |\textbf{Lago Eterno}|
Misión no encontrada.

=== PurisQuest ===
...
Opción: |\textbf{3}|

=== Perfil de Wilson ===
XP total: 150
Nivel: Aprendiz
Misiones completadas: ['Bosque Maldito']

=== PurisQuest ===
...
Opción: |\textbf{4}|
Hasta la próxima, Wilson. Nivel alcanzado: Aprendiz.
\end{minted}

\section*{{\small \faClipboardCheck\;} Rúbrica de evaluación}

\begin{table}[H]
    \centering
    \resizebox{0.6\textwidth}{!}{
    \begin{tabular}{p{11cm}r}
        \toprule
        \textbf{Criterios} & \textbf{Puntos} \\
        \midrule
        \multicolumn{2}{l}{\textit{Datos iniciales (5 pts)}}\\
        Diccionario \mintinline{python}{misiones} correctamente declarado & 3 \\
        Lista \mintinline{python}{completadas} vacía y \mintinline{python}{xp_acumulado} en \mintinline{python}{0} & 2 \\
        \midrule
        \multicolumn{2}{l}{\textit{Función \texttt{calcular\_nivel} (10 pts)}}\\
        Umbrales de nivel correctos & 6 \\
        Retorna el nombre del nivel como \mintinline{python}{str} & 4 \\
        \midrule
        \multicolumn{2}{l}{\textit{Función \texttt{mostrar\_misiones} (10 pts)}}\\
        Recorre el diccionario correctamente & 5 \\
        Imprime con el formato requerido & 5 \\
        \midrule
        \multicolumn{2}{l}{\textit{Función \texttt{completar\_mision} (20 pts)}}\\
        Normaliza el nombre & 4 \\
        Valida que la misión exista en el diccionario & 5 \\
        Valida que no haya sido completada previamente & 5 \\
        Agrega a \mintinline{python}{completadas} y acumula el XP correctamente & 4 \\
        Retorna el nuevo XP acumulado & 2 \\
        \midrule
        \multicolumn{2}{l}{\textit{Función \texttt{mostrar\_perfil} (10 pts)}}\\
        Imprime nombre, XP y lista de completadas & 5 \\
        Invoca \texttt{calcular\_nivel} para mostrar el nivel & 5 \\
        \midrule
        \multicolumn{2}{l}{\textit{Menú principal (15 pts)}}\\
        Solicita el nombre del jugador al inicio & 2 \\
        Ciclo \mintinline{python}{while} que se repite hasta la opción 4 & 5 \\
        Cada opción invoca la función correcta con los parámetros necesarios & 8 \\
        \midrule
        \textbf{Total} & \textbf{70} \\
        \bottomrule
    \end{tabular}
    }
    \caption{Rúbrica de evaluación de la Parte 1.}
    \label{tab:rubrica1}
\end{table}
